% =====================================================================
\section{Биты, байты и регистры}
% =====================================================================
\lessonintro
  {читать двоичные и шестнадцатеричные числа, управлять отдельными битами и выводить цифры на семисегментный индикатор через сдвиговый регистр}
  {урок 1 (цифровой выход), урок 10 (последовательная передача по такту)}
  {программа печатает биты числа в Serial Monitor; индикатор считает от 0 до 9, занимая всего три вывода платы}
  {микросхема 74HC595, семисегментный индикатор с общим катодом, восемь резисторов \SI{220}{\ohm}}

\subsection{Сначала: восемь выключателей в одном числе}

Внутри микроконтроллера нет «чисел» в привычном смысле --- есть ряды выключателей, каждый из которых
либо включён (1), либо выключен (0). Один выключатель --- это бит, восемь подряд --- байт. Число 91
и набор «выключен, включён, выключен, включён, включён, выключен, включён, включён» --- одно и то же.
Умея обращаться с отдельными битами, можно одним байтом зажечь нужные сегменты индикатора, прочитать
настройки датчика и понять любую таблицу регистров в документации.

\subsection{Двоичная и шестнадцатеричная запись}

У каждого бита свой «вес» --- степень двойки. Значение байта равно сумме весов включённых битов. Биты
нумеруют справа налево, начиная с нуля: бит 0 --- младший, бит 7 --- старший.

\begin{figure}[H]
\centering
\begin{tikzpicture}[x=1.25cm,y=1cm,font=\sffamily\small]
  \foreach \i/\b/\w in {0/0/128,1/1/64,2/0/32,3/1/16,4/1/8,5/0/4,6/1/2,7/1/1}{
    \pgfmathtruncatemacro{\n}{7-\i}
    \ifnum\b=1 \def\c{blkmem}\else\def\c{white}\fi
    \draw[fill=\c] (\i,0) rectangle ++(1,1);
    \node[font=\ttfamily\Large] at (\i+0.5,0.5) {\b};
    \node[font=\sffamily\scriptsize] at (\i+0.5,1.3) {бит \n};
    \node[font=\sffamily\scriptsize,text=accent] at (\i+0.5,-0.3) {\w};
  }
  \node[anchor=east,font=\sffamily\scriptsize,text=accent] at (-0.15,-0.3) {вес};
  \draw[decorate,decoration={brace,mirror,amplitude=4pt}] (0,-0.65) -- (3.98,-0.65) node[midway,below=5pt,font=\ttfamily]{0101 = 5};
  \draw[decorate,decoration={brace,mirror,amplitude=4pt}] (4.02,-0.65) -- (8,-0.65) node[midway,below=5pt,font=\ttfamily]{1011 = B};
  \node[anchor=west,align=left] at (8.4,0.5) {$64+16+8+2+1 = 91$\\[2pt] \texttt{0b01011011 = 0x5B}};
\end{tikzpicture}
\caption{Один байт в трёх записях: двоичной, десятичной и шестнадцатеричной. Каждая шестнадцатеричная
цифра заменяет ровно четыре бита}
\end{figure}

Шестнадцатеричная запись (hex) --- просто короткий способ записать биты: четыре бита превращаются в одну
цифру от 0 до F (A = 10, B = 11, \ldots, F = 15). В программе двоичное число пишут с приставкой
\code{0b}, шестнадцатеричное --- с \code{0x}. Адрес дисплея \code{0x3C} из урока 6 --- это \code{0b00111100},
или 60.

\subsection{Битовые операции}

\begin{center}\small
\begin{tabularx}{\textwidth}{@{}l l X@{}}
\toprule
Операция & Запись & Что делает \\
\midrule
И (AND) & \code{a \& b} & бит равен 1, только если он 1 в обоих числах; так \emph{проверяют} и \emph{сбрасывают} биты \\
ИЛИ (OR) & \code{a | b} & бит равен 1, если он 1 хотя бы в одном числе; так \emph{устанавливают} биты \\
Исключающее ИЛИ (XOR) & \code{a \textasciicircum{} b} & бит равен 1, если биты различаются; так \emph{переключают} биты \\
НЕ (NOT) & \code{\textasciitilde a} & меняет каждый бит на противоположный \\
Сдвиг & \code{a << n}, \code{a >> n} & сдвигает биты влево или вправо на $n$ позиций \\
\bottomrule
\end{tabularx}
\end{center}

Выражение \code{(1 << n)} --- это «маска»: число, в котором включён только бит $n$. С ней записываются
четыре приёма, которые встречаются в любой программе для микроконтроллера:

\begin{Verbatim}[fontsize=\small,bgcolor=codebg]
x |=  (1 << n);      // установить бит n
x &= ~(1 << n);      // сбросить бит n
x ^=  (1 << n);      // переключить бит n
if (x & (1 << n))    // проверить бит n
\end{Verbatim}

\codecap{Битовые операции: результат каждого шага виден в Serial Monitor}
\codeinput{l22_bits}

\begin{caution}
Не путайте \code{\&} и \code{\&\&}, \code{|} и \code{||}. Одиночные знаки работают с каждым битом числа,
двойные --- с условиями «истина/ложь». \code{if (a \& b)} и \code{if (a \&\& b)} --- разные проверки.
\end{caution}

\subsection{Регистры: память, которая управляет железом}

Регистр --- это ячейка памяти, биты которой напрямую связаны с аппаратурой. Запись числа в регистр
немедленно меняет поведение чипа. Функция \code{digitalWrite(4, HIGH)} в конце концов делает именно это:
записывает единицу в бит 4 регистра выходов. То же можно сделать вручную.

\codecap{Мигание записью в регистры GPIO напрямую}
\codeinput{l22_register}

В обычных программах пользуйтесь \code{digitalWrite()}: она понятнее и не зависит от модели чипа. Но
знать о регистрах нужно, чтобы читать документацию. Настройки любого датчика --- это тоже регистры,
только расположенные внутри датчика (урок 26).

\subsection{Сдвиговый регистр 74HC595}

Семисегментный индикатор --- это восемь светодиодов (семь сегментов и точка). Если подключать каждый
к своему выводу, на четыре цифры уйдёт вся плата. Микросхема 74HC595 решает задачу: она принимает
восемь битов по одному проводу и выставляет их на восьми своих выходах. Плате нужно всего три вывода.

\begin{figure}[H]
\centering
\begin{tikzpicture}[font=\sffamily\scriptsize]
  \node[blk,fill=blkcpu,minimum width=2.6cm,minimum height=2.6cm] (m) at (0,0) {ESP32-S3};
  \node[blk,fill=blkper,minimum width=2.8cm,minimum height=2.6cm] (r) at (6.4,0) {74HC595\\ сдвиговый\\ регистр};
  \draw[arr,accent] ($(m.east)+(0,0.8)$) -- node[above]{DATA (GPIO15) → DS} ($(r.west)+(0,0.8)$);
  \draw[arr,esp] ($(m.east)+(0,0)$) -- node[above]{CLOCK (GPIO16) → SHCP} ($(r.west)+(0,0)$);
  \draw[arr,okgreen] ($(m.east)+(0,-0.8)$) -- node[above]{LATCH (GPIO7) → STCP} ($(r.west)+(0,-0.8)$);
  \foreach \i/\l in {0/a,1/b,2/c,3/d,4/e,5/f,6/g,7/dp}{
    \draw[thick,gray] ($(r.east)+(0,1.05-\i*0.3)$) -- ++(1.5,0);
    \node[anchor=south,font=\sffamily\tiny] at ($(r.east)+(0.75,1.02-\i*0.3)$) {Q\i\ → \l};
  }
  \node[anchor=west,align=left] at ($(r.east)+(1.6,0)$) {через резисторы\\ \SI{220}{\ohm}\\ к сегментам};
  % индикатор
  \begin{scope}[shift={(12.8,-0.9)},line width=3pt,line cap=round,esp!80]
    \draw (0.15,1.8) -- (0.85,1.8);  \draw (1,1.65) -- (1,1.05);
    \draw (1,0.75) -- (1,0.15);      \draw (0.15,0) -- (0.85,0);
    \draw (0,0.75) -- (0,0.15);      \draw (0,1.65) -- (0,1.05);
    \draw (0.15,0.9) -- (0.85,0.9);  \fill (1.35,0) circle (1.6pt);
  \end{scope}
  \begin{scope}[shift={(12.8,-0.9)},font=\sffamily\tiny]
    \node at (0.5,2.0) {a}; \node at (1.2,1.35) {b}; \node at (1.2,0.45) {c}; \node at (0.5,-0.22) {d};
    \node at (-0.2,0.45) {e}; \node at (-0.2,1.35) {f}; \node at (0.5,1.08) {g}; \node at (1.6,0) {dp};
  \end{scope}
\end{tikzpicture}
\caption{Три вывода платы управляют восемью сегментами. Справа --- принятые названия сегментов индикатора}
\end{figure}

Работа идёт в три шага. Плата опускает LATCH. Затем восемь раз выставляет очередной бит на DATA и
«щёлкает» CLOCK --- биты один за другим вдвигаются в микросхему. Наконец LATCH поднимается, и все
восемь выходов обновляются одновременно. Всю передачу выполняет готовая функция \code{shiftOut()}.

Остальные выводы микросхемы: питание VCC (16) --- к 3V3, GND (8) --- к земле, $\overline{\text{OE}}$ (13) ---
к земле (выходы включены), $\overline{\text{MR}}$ (10) --- к 3V3 (сброс не активен). Общий катод индикатора ---
к GND.

\codecap{Счёт от 0 до 9 на семисегментном индикаторе}
\codeinput{l22_shift}

Каждое число в массиве \code{DIGITS} --- это рисунок цифры: единица в бите означает горящий сегмент.
Для цифры 2 горят a, b, d, e, g --- биты 0, 1, 3, 4, 6, то есть \code{0b01011011}. Это то самое число
с первого рисунка урока.

\begin{tip}
Микросхемы 74HC595 соединяют цепочкой: выход Q7S (вывод 9) первой подключают к входу DS второй, а CLOCK
и LATCH делают общими. Два вызова \code{shiftOut()} подряд --- и те же три вывода платы управляют уже
шестнадцатью выходами.
\end{tip}

\subsection{Самостоятельный опыт: счётчик нажатий}

Добавьте кнопку из урока 3: каждое нажатие увеличивает цифру на индикаторе, после 9 идёт 0. Затем
сделайте так, чтобы точка мигала раз в секунду, не мешая счёту. Подсказка: точку включает операция
\code{| 0b10000000}, а выключает --- её отсутствие.

\begin{trouble}
\begin{itemize}[leftmargin=*]
  \item \textbf{Индикатор не светится совсем} --- общий вывод индикатора не подключён к GND или вывод
        $\overline{\text{OE}}$ микросхемы «висит в воздухе».
  \item \textbf{Горят не те сегменты} --- перепутан порядок выходов Q0--Q7 или у вас индикатор с общим
        анодом: для него нужно инвертировать рисунок (\code{\textasciitilde pattern}) и подключить общий вывод к 3V3.
  \item \textbf{Цифры мелькают случайным образом} --- не подключён LATCH или перепутаны CLOCK и LATCH.
\end{itemize}
\end{trouble}

\begin{summary}
\begin{itemize}[leftmargin=*]
  \item Байт --- это восемь битов с весами 128\ldots1; hex-цифра заменяет четыре бита.
  \item Маска \code{(1 << n)} с операциями \code{|=}, \code{\&=\,\textasciitilde}, \code{\textasciicircum=} и \code{\&} устанавливает, сбрасывает, переключает и проверяет бит.
  \item Сдвиговый регистр превращает три вывода платы в восемь (и больше) выходов.
\end{itemize}
\end{summary}

\begin{selfcheck}
\begin{enumerate}
  \item Как перевести \code{0b00101101} в десятичную и шестнадцатеричную запись?
  \item Какой строкой кода сбросить бит 5 переменной \code{x}, не меняя остальные биты?
  \item Зачем сдвиговому регистру вывод LATCH, если данные уже переданы?
\end{enumerate}
\end{selfcheck}

\begin{tasks}
\begin{enumerate}
\item \lv{1} Переведите числа 13, 42 и 255 в двоичную и шестнадцатеричную запись и проверьте в программе.
\item \lv{1} Выведите состояние одного бита числа с помощью сдвига и маски.
\item \lv{2} Допишите в массив рисунки букв A, b, C, d, E, F и выведите шестнадцатеричный счёт на индикатор.
\item \lv{2} Включайте, выключайте и переключайте отдельный бит числа функциями-обёртками.
\item \lv{3} Подключите восемь светодиодов к 74HC595 и сделайте «бегущий огонь» одной операцией сдвига.
\item \lv{4} Соедините две 74HC595 цепочкой и выведите двузначное число на два индикатора.
\end{enumerate}
\end{tasks}
